\section{定数項のある漸化式}\label{節：定数項のある漸化式}
\subsection{不動点からのずれを追う}

まず
\[
 a_1=1,\qquad a_{n+1}=2a_n+1
\]
を考える. 項は$1,3,7,15,\ldots$と進む.
毎回1を足すため,\,元の数列は等比数列ではない.
しかし,\,各項に1を加えると
\[
 a_{n+1}+1=2(a_n+1)
\]
となる. $b_n=a_n+1$は初項2,\,公比2の等比数列だから
\[
 b_n=2^n,\qquad a_n=2^n-1.
\]

なぜ1を加えるのだろうか.
更新$f(x)=2x+1$で変わらない値を探すと,
$\alpha=2\alpha+1$から$\alpha=-1$となる.
このような$f(\alpha)=\alpha$を満たす値を\textbf{不動点}という.
先ほど追った$a_n+1$は,\,不動点からのずれ$a_n-\alpha$であった.

一般に
\[
 a_{n+1}=pa_n+q,\qquad p\ne1
\]
では,$\alpha=\frac{q}{1-p}$が不動点になる.
二つの式を引けば
\[
 a_{n+1}-\alpha=p(a_n-\alpha).
\]
従って
\[
 a_n=\alpha+p^{n-1}(a_1-\alpha).
\]
ただし$p=0$では,$n=1$は初項そのものを読み,
$n\ge2$では$a_n=q$とする.
零の零乗を計算する必要はない.

\NewFigFixed
原点を$\alpha$へ移せば,\,毎回の更新は$p$倍だけになる.
$p<0$なら向きが反転する.
$|p|<1$ならずれは零に近づき,\,数列は$\alpha$へ収束する.
$|p|>1$では,\,不動点から出発する場合を除いてずれの絶対値が増大する.
$p=-1$なら,\,不動点を挟んだ二つの値を交互に取る.
式変形は,\,数直線上の動きも説明していたのである.

\subsection{定数項も一つの成分にする}

この更新を行列で書くため,\,変化しない数1をもう一つの成分に加える.
\[
 \boldsymbol x_n=\begin{pmatrix}a_n\\1\end{pmatrix},
 \qquad
 A=\begin{pmatrix}p&q\\0&1\end{pmatrix}
\]
とすれば
\[
 \boldsymbol x_{n+1}=A\boldsymbol x_n.
\]
下の成分は常に1で,\,上の成分だけが元の規則に従う.
記録する成分を一つ増やすと,\,定数の足し算も行列の積に含められる.

$p\ne1$では
\[
 \begin{pmatrix}a_n\\1\end{pmatrix}
 =\begin{pmatrix}\alpha\\1\end{pmatrix}
 +(a_n-\alpha)\begin{pmatrix}1\\0\end{pmatrix}.
\]
前者は更新しても変わらず,\,後者の方向は$p$倍になる.
つまり,\,不動点とずれを分ける解法は,\,\textbf{変化の倍率が異なる二方向への分解}だった.

$p=1$の場合も行列の表示は使える.
\[
 \begin{pmatrix}1&q\\0&1\end{pmatrix}^{m}
 =\begin{pmatrix}1&mq\\0&1\end{pmatrix}
\]
なので,$a_n=a_1+(n-1)q$となる.
$q\ne0$では不動点がなく,\,ずれの倍率だけには直せない.
足し算が積み重なって$n$が現れることを覚えておこう.

\subsection{定数項がある連立漸化式}

定数ベクトル$\boldsymbol b$が加わる
\[
 \boldsymbol z_{n+1}=A\boldsymbol z_n+\boldsymbol b
\]
では,\,まず不動点$\boldsymbol z_*$を探す.
もし
\[
 \boldsymbol z_*=A\boldsymbol z_*+\boldsymbol b
\]
を満たせば,$\boldsymbol e_n=\boldsymbol z_n-\boldsymbol z_*$
によって
\[
 \boldsymbol e_{n+1}=A\boldsymbol e_n
\]
に戻る. 一つの数を更新する場合と同じように,不動点からのずれを追っている.

例えば
\[
 \begin{aligned}
 x_{n+1}&=2x_n+y_n+1,\\
 y_{n+1}&=x_n+2y_n+1
 \end{aligned}
\]
の不動点は$(0,-1)$である.
$u_n=x_n,\,v_n=y_n+1$と置けば,\,定数項のない冒頭の連立式に戻る.
ただし不動点が存在するとは限らず,\,複数ある場合もある.
その有無は$(I-A)\boldsymbol z_*=\boldsymbol b$という連立方程式で確かめる.

不動点が見つからなくても,\,一定成分1を加えれば
\[
 \begin{pmatrix}\boldsymbol z_{n+1}\\1\end{pmatrix}
 =
 \begin{pmatrix}A&\boldsymbol b\\0&1\end{pmatrix}
 \begin{pmatrix}\boldsymbol z_n\\1\end{pmatrix}
\]
と書ける. 右下の0は,\,必要な個数の零を横に並べた部分を表す.
次章では,隣り合う項を成分に選び,一つの数列を連立漸化式として調べる.


\begin{bookpoint}
一つの数でも複数の数でも,不動点を引けば定数項が消える.
不動点を持たない場合にも,一定成分1を加えれば行列で更新できる.
\end{bookpoint}
